% from main.tex
%%%%%%%%%%  scan idx 254  |  folio 241  %%%%%%%%%%
\emph{Preuve.} Conséquence de (2.2), (2.5) et (1.12).

\textbf{Scholie 2.6.} Pour démontrer la conjecture (2.1) dans le cas général, il faudra sans doute mettre en œuvre de nouvelles techniques. On peut espérer que celles-ci aideront un jour à prouver la conjecture analogue en Géométrie Analytique, pour un morphisme propre d'espaces analytiques au-dessus d'un topos localement annelé en \(\mathbb C\)-algèbres convenablement topologisées, conjecture qui n'est « démontrée » pour l'instant qu'au-dessus d'un espace analytique complexe (théorème de Grauert [1]).\footnote{(ajouté en Avril 1971) voir un article en préparation de R. Kiehl et J.L. Verdier, où le théorème de finitude 1 est démontré par des techniques différentes, plus simples, et sous des hypothèses plus générales.}

\section*{3. Tor-dimension finie relative}

\textbf{Définition 3.1.} Soient \(f:(X,\mathcal O_X)\longrightarrow(Y,\mathcal O_Y)\) un morphisme de topos annelés, \(E\in\Ob D^{-}(X)\), \([m,n]\) un intervalle de \(\mathbb Z\). On dit que \(E\) est d'amplitude plate relativement à \(f\) (ou \(Y\), s'il n'y a pas de confusion possible) contenue dans \([m,n]\), et l'on écrit
\[
  \toramp_f(E)\subset[m,n]\qquad(\text{ou }\toramp_Y(E)\subset[m,n])
\]
si \(E\), considéré comme complexe de \(f^{-1}(\mathcal O_Y)\)-Modules, est d'amplitude plate contenue dans \([m,n]\) (I 5.2). On dit que \(f\) est d'amplitude plate contenue dans \([-a,0]\) \((a\geq0)\), ou de tor-dimension \(\leq a\), si \(\mathcal O_X\) est d'amplitude plate rel. à \(f\), contenue dans \([-a,0]\). On dit que \(E\) est d'amplitude plate finie rel. à \(f\) s'il existe un intervalle \([m,n]\) tel que l'on ait \(\toramp_f(E)\subset[m,n]\).

\textbf{Notation 3.2.} On notera \(D(f)_{\mathrm{torf}}\), ou \(D(X)_{f\text{-}\torf}\), ou \(D(X)_{Y\text{-}\torf}\) s'il n'y a pas de confusion possible, la sous-catégorie pleine de \(D^{-}(X)\) formée des complexes d'amplitude plate finie relativement à \(f\). En vertu de (I 5.3), c'est une sous-catégorie triangulée de \(D^b(X)\).

